\chapter{Cardiac Stress Test}

\section*{What Is This Test?}
A cardiac stress test evaluates how the heart responds when it works harder than it does at rest. In an exercise stress test, the patient walks on a treadmill or pedals a stationary bicycle while heart rhythm, electrocardiogram (ECG), blood pressure, symptoms, and exercise capacity are monitored. If the patient cannot exercise adequately, medication may be used to increase heart workload or change coronary blood flow. Stress testing may be combined with echocardiography or nuclear myocardial-perfusion imaging to provide information beyond a standard exercise ECG.

\section*{Alternative Names}
\begin{itemize}
  \item Exercise Stress Test
  \item Treadmill Test
  \item Exercise ECG
  \item Cardiac Stress Test
  \item Pharmacologic Stress Test
  \item Stress Echocardiogram
  \item Nuclear Stress Test
  \item Myocardial Perfusion Stress Test
\end{itemize}
\section*{Principle}
Exercise or a pharmacologic agent increases the heart's oxygen demand or alters blood flow. Continuous ECG and serial blood-pressure measurements show the electrical and hemodynamic response. Exercise-induced ischemia may cause characteristic ECG changes, symptoms, abnormal blood-pressure response, or regional wall-motion or perfusion abnormalities on imaging. A stress test estimates the likelihood of certain cardiac conditions; it does not by itself prove or exclude every form of coronary artery disease.

\section*{Why This Test Is Ordered}
\begin{itemize}
  \item Evaluation of exertional chest discomfort, shortness of breath, unusual fatigue, or lightheadedness when a cardiac cause is suspected
  \item Assessment of suspected coronary artery disease in an appropriate clinical setting
  \item Evaluation of exercise-induced arrhythmias or an abnormal response to exertion
  \item Measurement of functional exercise capacity and blood-pressure response
  \item Assessment of response to treatment or cardiac rehabilitation when clinically indicated
  \item Risk assessment or preoperative evaluation when the expected result would change management
\end{itemize}

\section*{Primary vs Secondary Test}
The exercise ECG is a primary functional stress test when the patient can exercise adequately and the resting ECG is interpretable. Stress echocardiography or nuclear perfusion imaging is a secondary or alternative approach when more diagnostic information is needed, when the resting ECG limits interpretation, or when the patient cannot exercise sufficiently. The choice depends on symptoms, cardiovascular risk, baseline ECG, ability to exercise, kidney function, pregnancy status, and local expertise.

\section*{How This Test Is Performed}
ECG electrodes are placed on the chest and a blood-pressure cuff is applied. After a resting ECG and baseline measurements, treadmill speed and incline or bicycle resistance increase in stages. The patient reports chest discomfort, shortness of breath, dizziness, leg fatigue, or other symptoms. The test stops when the target workload is reached, symptoms or ECG changes require stopping, blood pressure responds abnormally, or the patient asks to stop. A pharmacologic study uses an intravenous medicine and may include an imaging tracer or ultrasound images. Recovery is monitored after exercise or medication.

\section*{What Preparation Is Needed}
Follow the testing center's instructions. Wear comfortable walking shoes and clothing, avoid a heavy meal beforehand, and do not smoke immediately before the test. Ask the ordering clinician whether beta-blockers, calcium-channel blockers, nitrates, or other medicines should be withheld; do not stop prescribed medicines without instructions. Tell the team about pregnancy, recent illness, asthma, diabetes, mobility limitations, allergies, and any implanted device. Caffeine restrictions may apply to some pharmacologic nuclear tests.

\section*{Contraindications and Precautions}
\begin{itemize}
  \item The test requires individualized medical review and should not be performed during an acute myocardial infarction or unstable angina, uncontrolled symptomatic arrhythmia, severe symptomatic aortic stenosis, decompensated heart failure, acute pulmonary embolism, acute myocarditis or pericarditis, or other unstable illness. Exercise may be unsuitable when the patient cannot walk or pedal safely. A normal or abnormal result must be interpreted with the pretest likelihood, test quality, imaging findings, and symptoms; it should not be used as a stand-alone diagnosis.
\end{itemize}
\section*{Risks}
Temporary breathlessness, fatigue, muscle discomfort, palpitations, dizziness, or chest discomfort may occur. Rare but serious complications include sustained arrhythmia, severe blood-pressure changes, myocardial infarction, or cardiac arrest. The test is performed with trained staff and emergency equipment available. Pharmacologic tests can cause medication-specific symptoms such as flushing, headache, nausea, wheezing, or transient blood-pressure changes.

\section*{What Happens After the Test}
The patient is monitored during recovery until heart rate, blood pressure, and symptoms return toward baseline. Results may be available immediately for an exercise ECG, while imaging interpretation may take longer. An abnormal, equivocal, or technically limited study may lead to additional imaging, coronary CT angiography, ambulatory ECG monitoring, echocardiography, or cardiology review. Emergency assessment is required for persistent or severe chest pain, fainting, or breathing difficulty after the test.

\section*{Understanding Results}

\subsection*{Negative or Reassuring}
The patient reaches an adequate workload without concerning symptoms, significant ischemic ECG changes, abnormal blood-pressure response, or imaging evidence of inducible ischemia. The result lowers the likelihood of significant flow-limiting coronary disease under the conditions tested but does not exclude plaque, microvascular disease, or future cardiac events.

\subsection*{Positive or Abnormal}
The study shows symptoms, ECG changes, an abnormal blood-pressure or heart-rate response, an exercise-induced arrhythmia, or a new regional wall-motion or perfusion abnormality. The result may indicate inducible myocardial ischemia or another cardiac problem and requires clinician interpretation. False-positive and false-negative results occur, particularly when exercise is inadequate, the resting ECG is abnormal, or the patient's pretest probability is very low or very high.

\section*{Follow-up Tests}
\begin{itemize}
  \item \textbf{Stress echocardiography or nuclear myocardial-perfusion imaging:} To clarify an abnormal or nondiagnostic exercise ECG or assess the location and extent of inducible ischemia.
  \item \textbf{Coronary CT angiography or invasive coronary angiography:} Considered when symptoms and risk indicate an anatomic assessment is needed.
  \item \textbf{Resting echocardiogram:} To evaluate ventricular function, valves, and structural heart disease.
  \item \textbf{Ambulatory ECG monitoring:} To investigate intermittent palpitations, syncope, or exertional arrhythmias.
  \item \textbf{Cardiovascular risk assessment:} Blood pressure, lipid profile, glucose or HbA1c, smoking status, and family history help guide overall risk reduction.
\end{itemize}

\section*{References}
\begin{enumerate}
  \item American Heart Association. Exercise Stress Test. Last reviewed February 24, 2025.
  \item Gulati M, Levy PD, Mukherjee D, et al. 2021 AHA/ACC/ASE/CHEST/SAEM/SCCT/SCMR guideline for the evaluation and diagnosis of chest pain. \textit{Circulation}. 2021;144:e368--e454.
  \item Mayo Clinic. Stress test. Updated 2024.
\end{enumerate}

\clearpage
\section*{Regulatory Status and Authorization Date}
This chapter describes a test category, not a specific commercial product. FDA, CDSCO, EU IVDR, and MHRA approval or registration dates therefore cannot be assigned without the assay manufacturer, product name, instrument, and intended use. For laboratory-developed tests, an individual agency approval date may not exist. See the Preface for the interpretation of regulatory status.

\clearpage
